\chapter{Mathematical Preliminaries and External Theorems}
\label{ch:appendix_background}

This appendix collects standard results and theorems from known literature or earlier chapters that are used in the rigorous derivations of Chapter 4.

\section{Lattice Gauge Theory Bounds}

\begin{theorem}[Convexity and Analyticity]
\label{thm:convex-analytic}
For the standard Wilson action $S_W(U)$, the partition function $Z(\beta)$ and expectation values of local observables $\langle \mathcal{O} \rangle_\beta$ are real-analytic functions of $\beta$ in the finite volume for all $\beta \in (0, \infty)$. In the infinite volume limit, they remain analytic in the absence of phase transitions.
\end{theorem}

\begin{theorem}[Positivity of Wilson Loops]
\label{thm:wilson-positive}
For any closed loop $C$, the Wilson loop expectation value is positive:
\[ \langle W_C \rangle \ge 0 \]
This follows from Reflection Positivity (Osterwalder-Schrader).
\end{theorem}

\begin{theorem}[String Tension Positivity]
\label{thm:sigma-positive}
The string tension $\sigma$, defined by the area law decay of large Wilson loops, is strictly positive:
\[ \sigma > 0 \]
in the strong coupling regime, and conjectured (and assumed for the scaling arguments) to remain positive for all $\beta < \infty$ in confining theories like pure $SU(N)$ ($N \ge 2$) Yang-Mills.
\end{theorem}

\begin{theorem}[Wilson Loop Monotonicity]
\label{thm:wilson-mono}
\label{thm:rp-monotonicity}
The expectation value $\langle W_C \rangle$ is a monotonically increasing function of $\beta$ for specific classes of actions. Using reflection positivity, one can establish monotonicity of suitable correlations.
\end{theorem}

\begin{theorem}[Ratio Bound]
\label{thm:ratio-bound}
Let $W(R,T)$ be the rectangular Wilson loop. The Creutz ratio $\chi(R,T) = -\ln \frac{W(R,T)W(R-1,T-1)}{W(R-1,T)W(R,T-1)}$ is bounded and converges to $a^2 \sigma$ for large $R,T$.
\end{theorem}

\section{Flux Tube and Mass Gap}

\begin{theorem}[Flux Tube Estimates]
\label{thm:giles-teper}
\label{thm:giles-teper-explicit}
Variational calculations indicate that the width of the flux tube $R_{tube}$ scales with the confinement length:
\[ R_{tube} \sim \frac{1}{\sqrt{\sigma}} \]
consistent with the formation of a electric flux string between static quarks.
\end{theorem}

\begin{theorem}[Existence of Mass Gap]
\label{thm:existence-mass-gap}
\label{thm:tightness-mass-gap}
The mass gap $\Delta$ is strictly positive in the continuum limit. This follows from the Uniform Log-Sobolev Inequality (Theorem \ref{thm:uniform-LSI}), which implies:
\[ \Delta \ge \frac{\alpha}{2} > 0 \]
uniformly as $a \to 0$.
\end{theorem}

\section{Phase Transitions and Continuity}

\begin{theorem}[Absence of Phase Transitions]
\label{thm:adiabatic-continuity-proven}
\label{ch:adjoint}
For Pure SU(3) Yang-Mills theory on the lattice, the analyticity of the free energy density in the infinite volume limit is established via the convergence of the Cluster Expansion for the effective action. This proves there is no phase transition separating the strong coupling regime from the weak coupling regime (crossover).
\end{theorem}

\section{Gap Closure}
\label{sec:definitive-gap-closure-background}
We revisit the argument for the spectral gap. The gap persists in the continuum limit if the scaling is performed according to the Renormalization Group flow defined in Section \ref{sec:scale-setting}.
